%! Trigonometric Equations and Inequalities — Unit 7, Lesson 7.3 — Guided Notes
\documentclass[10pt]{article}
\usepackage{precalculus-article}
\usepackage{precalculus-boxes}
\newcommand{\TallMath}[1]{$\displaystyle #1\rule[-1.4em]{0pt}{3.2em}$}

\begin{document}

\pageheader{Unit 7, Lesson 7.3}{Guided Notes}
\namedateperiod

\vspace{0.08in}

%----------------------------------------------------------------------
% Objectives
%----------------------------------------------------------------------
\begin{objectivebox}
\textbf{I will be able to\ldots}
\begin{enumerate}[leftmargin=1.5em, itemsep=4pt]
  \item \textbf{LO 3.10.A} --- Solve a trigonometric equation by isolating
    the trig function and using inverse trig to find all solutions.
    \writeline
  \item \textbf{LO 3.10.A} --- Write the general solution of a trig equation
    using periodicity, and apply a domain restriction when context requires it.
    \writeline
\end{enumerate}
\end{objectivebox}

\vspace{0.06in}

%----------------------------------------------------------------------
% Vocabulary
%----------------------------------------------------------------------
\begin{vocabbox}
\begin{description}[leftmargin=0pt, itemsep=5pt]
  \item[\termblanklong{trigonometric equation}]
    An equation in which the unknown appears inside a
    \underline{\hspace{4cm}} function.

  \item[\termblanklong{principal solution}]
    The solution obtained by applying an \underline{\hspace{3.5cm}} trig function
    directly; limited to the \underline{\hspace{3cm}} of that inverse function.

  \item[\termblanklong{general solution}]
    \textbf{All} solutions, expressed by adding integer multiples of the
    \underline{\hspace{3cm}} to the principal value(s).
    Example: $\theta = \tfrac{\pi}{6} + $ \underline{\hspace{1.5cm}},
    $k \in \mathbb{Z}$.

  \item[\termblanklong{domain restriction}]
    A limit on the variable imposed by the \underline{\hspace{4cm}} of the
    problem (e.g.\ $0 \le t \le 24$ for a 24-hour model).
\end{description}
\end{vocabbox}

\vspace{0.06in}

%----------------------------------------------------------------------
% LO 3.10.A — Section 1: Infinitely Many Solutions
%----------------------------------------------------------------------
\begin{notesbox}{LO 3.10.A --- EK 3.10.A.2: Infinitely Many Solutions}

\textbf{Key Idea:} Because trig functions are \underline{\hspace{3cm}},
a trig equation typically has \underline{\hspace{3cm}} many solutions when
the domain is all real numbers.

\vspace{6pt}
\textbf{Example:} Solve $\sin\theta = \dfrac{1}{2}$.

\vspace{4pt}
\textbf{Step 1 --- Principal value:}
$\theta = \arcsin\!\bigl(\tfrac{1}{2}\bigr) = $ \underline{\hspace{2cm}}

\vspace{4pt}
\textbf{Step 2 --- Second solution in $[0, 2\pi)$:}
Sine is positive in Quadrants \underline{\hspace{1cm}} and \underline{\hspace{1cm}}.
Reflect: $\theta = \pi - \underline{\hspace{1.5cm}} = $ \underline{\hspace{2cm}}

\vspace{4pt}
\textbf{Step 3 --- General solution:}
Add multiples of the period $2\pi$ to each solution.

\vspace{4pt}
\renewcommand{\arraystretch}{1.8}
\begin{tabular}{@{}l@{\hspace{2cm}}l@{}}
Family 1: $\theta = $ \underline{\hspace{3cm}} & Family 2: $\theta = $ \underline{\hspace{3cm}}
\end{tabular}

\vspace{4pt}
where $k$ is any \underline{\hspace{3cm}}.

\vspace{6pt}
\textbf{Check:} List 3 specific solutions from each family:

Family 1 ($k = 0, 1, 2$): \writeline

Family 2 ($k = 0, 1, -1$): \writeline

\end{notesbox}

\vspace{0.06in}

%----------------------------------------------------------------------
% Section 2: Inverse Trig Procedure
%----------------------------------------------------------------------
\begin{notesbox}{LO 3.10.A --- EK 3.10.A.1: The Four-Step Procedure}

\textbf{Algorithm for solving a trig equation:}
\begin{enumerate}[leftmargin=1.5em, itemsep=4pt]
  \item \textbf{Isolate} the trig function on one side of the equation.
  \item \textbf{Apply} the inverse trig function to find the principal value.
  \item \textbf{Find} the second solution in $[0, 2\pi)$ using reference angles
    and the appropriate quadrant(s).
  \item \textbf{Write} the general solution by adding $+2\pi k$ (or $+\pi k$
    for tangent) to each solution.
\end{enumerate}

\vspace{6pt}
\textbf{Example:} Solve $2\cos\theta - 1 = 0$.

\vspace{4pt}
Step 1 --- Isolate: $\cos\theta = $ \underline{\hspace{2.5cm}}

\vspace{4pt}
Step 2 --- Principal value: $\theta = \arccos\!\bigl(\underline{\hspace{1cm}}\bigr) = $ \underline{\hspace{2cm}}

\vspace{4pt}
Step 3 --- Cosine is positive in Q\underline{\hspace{0.5cm}} and Q\underline{\hspace{0.5cm}}.
Second solution: $\theta = $ \underline{\hspace{3cm}}

\vspace{4pt}
Step 4 --- General solution:

$\theta = $ \underline{\hspace{3.5cm}} or $\theta = $ \underline{\hspace{3.5cm}},
$k \in \mathbb{Z}$

\vspace{6pt}
\textbf{Solutions in $[0, 2\pi)$:}
\writeline

\end{notesbox}

\vspace{0.06in}

%----------------------------------------------------------------------
% Section 3: Contextual Domain
%----------------------------------------------------------------------
\begin{notesbox}{LO 3.10.A --- EK 3.10.A.3: Contextual Domain Restrictions}

\textbf{Context:} The height of water at a harbor (in feet) is modeled by
$h(t) = 3\sin\!\bigl(\tfrac{\pi}{6}t\bigr) + 5$, where $t$ is hours after
midnight, $0 \le t \le 24$.

\vspace{6pt}
\textbf{Question:} When is $h(t) \ge 7$? (When can a ship requiring 7 ft enter?)

\vspace{4pt}
Step 1 --- Rewrite: $3\sin\!\bigl(\tfrac{\pi}{6}t\bigr) + 5 \ge 7$
\hspace{1em}$\Rightarrow$\hspace{1em}
$\sin\!\bigl(\tfrac{\pi}{6}t\bigr) \ge $ \underline{\hspace{2cm}}

\vspace{4pt}
Step 2 --- Let $u = \tfrac{\pi}{6}t$. Solve $\sin u = \tfrac{2}{3}$ for the boundary:

\vspace{3pt}
$u_1 = \arcsin\!\bigl(\tfrac{2}{3}\bigr) \approx $ \underline{\hspace{2cm}} rad
\hspace{1cm}
$u_2 = \pi - \arcsin\!\bigl(\tfrac{2}{3}\bigr) \approx $ \underline{\hspace{2cm}} rad

\vspace{4pt}
Step 3 --- Solve for $t$: $t = \dfrac{6}{\pi} \cdot u$

\vspace{3pt}
$t_1 \approx $ \underline{\hspace{2cm}} hr
\hspace{1.5cm}
$t_2 \approx $ \underline{\hspace{2cm}} hr

\vspace{4pt}
Step 4 --- The inequality $h(t) \ge 7$ holds on the interval
$\underline{\hspace{1.5cm}} \le t \le \underline{\hspace{1.5cm}}$.

\vspace{4pt}
\textbf{Interpret:} The ship can enter between approximately
\underline{\hspace{1.5cm}} and \underline{\hspace{1.5cm}} hours after midnight.

\vspace{4pt}
\textbf{Why do we restrict to $[0, 24]$?} \writeline

\end{notesbox}

\vspace{0.06in}

%----------------------------------------------------------------------
% Practice
%----------------------------------------------------------------------
\begin{practicebox}

\textbf{Guided Practice 1.}\quad Solve $\sqrt{2}\,\sin\theta - 1 = 0$ for
all $\theta \in [0, 2\pi)$. Show the four steps.
\writelines{4}

\vspace{0.08in}
\textbf{Guided Practice 2.}\quad Write the general solution to $\cos\theta = -\dfrac{\sqrt{3}}{2}$.

\vspace{4pt}
$\theta = $ \blank{4.5cm} or $\theta = $ \blank{4.5cm}, $k \in \mathbb{Z}$

\vspace{0.08in}
\textbf{Guided Practice 3.}\quad A temperature model gives $T(t) = 10\sin\!\bigl(\tfrac{\pi}{12}t\bigr) + 65$
for $t \in [0, 24]$ (hours). For what hours is $T(t) \ge 70$?

\vspace{4pt}
Set up the inequality and isolate the trig function:
\writelines{2}

Solve for the boundary values of $t$:
\writelines{2}

Answer: $T(t) \ge 70$ on \blank{5cm}.

\end{practicebox}

\end{document}
